% !PRINTED 504
\weberpage{504}

For the singular value all the quantities (12) now become equal
to one another and equal to $u$; $\sigma$ is therefore a rational function of
$u$, and with this the desired proof is completed.

As noteworthy examples we choose the cases of § 125:
\[
\begin{aligned}
\gamma_3\left(\frac{-1+\sqrt{-3}}{2}\right) &= 24\sqrt{-3}, \\
\gamma_3\left(\frac{-1+\sqrt{-7}}{2}\right) &= 27\sqrt{-7}, \\
\gamma_3\left(\frac{-1+\sqrt{-11}}{2}\right) &= 7.\,8\sqrt{-11}, \\
\gamma_3\left(\frac{-1+\sqrt{-19}}{2}\right) &= 8.\,27\sqrt{-19}, \\
\gamma_3\left(\frac{-1+\sqrt{-27}}{2}\right) &= 8.\,11.\,23\sqrt{-3}, \\
\gamma_3\left(\frac{-1+\sqrt{-43}}{2}\right) &= 8.\,7.\,81\sqrt{-43}, \\
\gamma_3\left(\frac{-1+\sqrt{-67}}{2}\right) &= 7.\,8.\,9.\,31\sqrt{-67}, \\ % [sic] printed 9; the arithmetic would require 27
\gamma_3\left(\frac{-1+\sqrt{-163}}{2}\right) &= 8.\,27.\,7.\,11.\,19.\,127\sqrt{-163}.
\end{aligned}
\]
Here too the decomposability of the rational factors into comparatively small primes is striking. We give here a few
further examples in which it is not $\sqrt{D}$ but $\sqrt{-D}$ that occurs.

In § 128 we found:
\[
f_1(\sqrt{-2})^4 = 2, \quad f(\sqrt{-3})^3 = 2, \quad f(\sqrt{-7})^2 = 2.
\]
To these we apply the formulae (§ 129):
\[
f_1^8 - f_2^8 = \frac{\sqrt{f^{24} - 64}}{f^4},
\]
\[
f^8 + f_2^8 = \frac{\sqrt{f_1^{24} + 64}}{f_1^4},
\]
and find
\[
\begin{aligned}
f(\sqrt{-2})^8 + f_2(\sqrt{-2})^8 &= 4\sqrt{2}, \\
f_1(\sqrt{-3})^8 - f_2(\sqrt{-3})^8 &= \frac{4\sqrt{3}}{\sqrt[3]{2}}, \\
f_1(\sqrt{-7})^8 - f_2(\sqrt{-7})^8 &= 6\sqrt{7}.
\end{aligned}
\]
